\section{Learning Objectives}

After completing this chapter, readers will be able to:
\begin{itemize}
    \item Apply Bayesian inference to A/B testing
    \item Construct and interpret posterior distributions
    \item Use conjugate priors for analytical solutions
    \item Calculate probability of superiority and expected loss
    \item Implement Bayesian A/B tests in Python and R
    \item Understand advantages and limitations of Bayesian approaches
\end{itemize}

\section{Bayesian vs. Frequentist Paradigms}

The Bayesian framework treats parameters as random variables and updates beliefs using observed data.

\subsection{Key Philosophical Differences}

\begin{center}
\begin{tabular}{lll}
\toprule
\textbf{Aspect} & \textbf{Frequentist} & \textbf{Bayesian} \\
\midrule
Parameters & Fixed constants & Random variables \\
Probability & Long-run frequency & Degree of belief \\
Inference & Confidence intervals & Credible intervals \\
Prior info & Not incorporated & Explicitly modeled \\
Interpretation & Indirect (via procedures) & Direct (probability) \\
Sequential testing & Problematic & Natural \\
\bottomrule
\end{tabular}
\end{center}

\subsection{Bayes' Theorem for A/B Testing}

\begin{equation}
p(\theta | D) = \frac{p(D | \theta) \cdot p(\theta)}{p(D)} \propto p(D | \theta) \cdot p(\theta)
\end{equation}

where:
\begin{itemize}
    \item $\theta$: Parameter of interest (e.g., conversion rate)
    \item $D$: Observed data
    \item $p(\theta)$: Prior distribution (before seeing data)
    \item $p(D | \theta)$: Likelihood (probability of data given parameter)
    \item $p(\theta | D)$: Posterior distribution (after seeing data)
\end{itemize}

The posterior combines prior beliefs with observed evidence.

\section{Conjugate Priors for Common Distributions}

Conjugate priors yield posterior distributions in the same family, enabling analytical solutions.

\subsection{Beta-Binomial Model for Conversion Rates}

For binary outcomes (conversions), the natural model is:

\textbf{Likelihood:}
\begin{equation}
X | p \sim \text{Binomial}(n, p)
\end{equation}

\textbf{Prior:}
\begin{equation}
p \sim \text{Beta}(\alpha, \beta)
\end{equation}

The Beta distribution has PDF:
\begin{equation}
f(p; \alpha, \beta) = \frac{\Gamma(\alpha + \beta)}{\Gamma(\alpha)\Gamma(\beta)} p^{\alpha-1}(1-p)^{\beta-1}, \quad p \in (0, 1)
\end{equation}

Properties:
\begin{itemize}
    \item $\mathbb{E}[p] = \frac{\alpha}{\alpha + \beta}$
    \item $\text{Mode}[p] = \frac{\alpha - 1}{\alpha + \beta - 2}$ for $\alpha, \beta > 1$
    \item $\text{Var}(p) = \frac{\alpha\beta}{(\alpha+\beta)^2(\alpha+\beta+1)}$
\end{itemize}

\textbf{Posterior:}

Given $x$ conversions from $n$ trials:

\begin{theorem}[Beta-Binomial Conjugacy]
If $p \sim \text{Beta}(\alpha, \beta)$ and $X | p \sim \text{Binomial}(n, p)$, then:
\begin{equation}
p | X=x \sim \text{Beta}(\alpha + x, \beta + n - x)
\end{equation}
\end{theorem}

\textbf{Interpretation:}
\begin{itemize}
    \item Prior $\text{Beta}(\alpha, \beta)$ represents $\alpha - 1$ prior successes and $\beta - 1$ prior failures
    \item Posterior adds observed data: $x$ successes and $n - x$ failures
    \item Total: $\alpha + x - 1$ successes, $\beta + n - x - 1$ failures
\end{itemize}

\begin{example}[Uninformative Prior]
Use $\text{Beta}(1, 1)$ (uniform on $[0,1]$) when no prior information exists.

After observing 80 conversions from 1000 users:
\begin{equation}
p | D \sim \text{Beta}(1 + 80, 1 + 920) = \text{Beta}(81, 921)
\end{equation}

Posterior mean: $\mathbb{E}[p | D] = \frac{81}{81 + 921} = 0.0808$

95\% credible interval: $[0.0649, 0.0992]$ (using beta quantiles)
\end{example}

\lstset{language=Python, caption=Bayesian Inference for Conversion Rate}
\begin{lstlisting}
import numpy as np
from scipy.stats import beta
import matplotlib.pyplot as plt

def bayesian_conversion_inference(conversions, trials,
                                   prior_alpha=1, prior_beta=1):
    """
    Bayesian inference for conversion rate using Beta-Binomial model.

    Parameters:
    -----------
    conversions : int
        Number of conversions
    trials : int
        Number of trials
    prior_alpha, prior_beta : float
        Beta prior parameters

    Returns:
    --------
    dict : Posterior statistics
    """
    # Posterior parameters
    post_alpha = prior_alpha + conversions
    post_beta = prior_beta + trials - conversions

    # Posterior statistics
    post_mean = post_alpha / (post_alpha + post_beta)
    post_mode = (post_alpha - 1) / (post_alpha + post_beta - 2) \
                if post_alpha > 1 and post_beta > 1 else None
    post_var = (post_alpha * post_beta) / \
               ((post_alpha + post_beta)**2 * (post_alpha + post_beta + 1))

    # Credible intervals
    ci_95 = beta.ppf([0.025, 0.975], post_alpha, post_beta)
    ci_90 = beta.ppf([0.05, 0.95], post_alpha, post_beta)

    return {
        'posterior_alpha': post_alpha,
        'posterior_beta': post_beta,
        'posterior_mean': post_mean,
        'posterior_mode': post_mode,
        'posterior_std': np.sqrt(post_var),
        'ci_95': ci_95,
        'ci_90': ci_90,
        'observed_rate': conversions / trials
    }

# Example
result = bayesian_conversion_inference(conversions=80, trials=1000)

print("Posterior Beta({}, {})".format(result['posterior_alpha'],
                                      result['posterior_beta']))
print(f"Posterior mean: {result['posterior_mean']:.4f}")
print(f"Posterior std: {result['posterior_std']:.4f}")
print(f"95% Credible Interval: [{result['ci_95'][0]:.4f}, "
      f"{result['ci_95'][1]:.4f}]")

# Visualize posterior
p_values = np.linspace(0.05, 0.12, 1000)
posterior_pdf = beta.pdf(p_values, result['posterior_alpha'],
                         result['posterior_beta'])

plt.figure(figsize=(10, 6))
plt.plot(p_values, posterior_pdf, linewidth=2, label='Posterior')
plt.axvline(result['posterior_mean'], color='red', linestyle='--',
            label=f"Mean = {result['posterior_mean']:.4f}")
plt.fill_between(p_values, 0, posterior_pdf,
                where=(p_values >= result['ci_95'][0]) &
                      (p_values <= result['ci_95'][1]),
                alpha=0.3, label='95% Credible Interval')
plt.xlabel('Conversion Rate')
plt.ylabel('Density')
plt.title('Posterior Distribution for Conversion Rate')
plt.legend()
plt.grid(alpha=0.3)
plt.tight_layout()
plt.savefig('bayesian_posterior.pdf')
plt.show()
\end{lstlisting}

\subsection{Normal-Normal Model for Continuous Metrics}

For continuous outcomes with known variance:

\textbf{Likelihood:}
\begin{equation}
Y_i | \mu \sim \mathcal{N}(\mu, \sigma^2)
\end{equation}

\textbf{Prior:}
\begin{equation}
\mu \sim \mathcal{N}(\mu_0, \tau^2)
\end{equation}

\textbf{Posterior:}

Given data $\mathbf{y} = (y_1, \ldots, y_n)$ with sample mean $\bar{y}$:

\begin{theorem}[Normal-Normal Conjugacy]
\begin{equation}
\mu | \mathbf{y} \sim \mathcal{N}\left(\frac{\frac{\mu_0}{\tau^2} + \frac{n\bar{y}}{\sigma^2}}{\frac{1}{\tau^2} + \frac{n}{\sigma^2}}, \left(\frac{1}{\tau^2} + \frac{n}{\sigma^2}\right)^{-1}\right)
\end{equation}
\end{theorem}

The posterior mean is a precision-weighted average of prior mean and sample mean.

\section{Bayesian A/B Testing}

For comparing two variants, we need the posterior distribution of the difference or ratio.

\subsection{Comparing Two Conversion Rates}

Given:
\begin{align}
p_A &\sim \text{Beta}(\alpha_A, \beta_A) \\
p_B &\sim \text{Beta}(\alpha_B, \beta_B)
\end{align}

We want to know:
\begin{equation}
P(p_B > p_A | D) = \int_0^1 \int_0^{p_B} f(p_A | D) f(p_B | D) \, dp_A \, dp_B
\end{equation}

No closed form, but easily computed via:
\begin{enumerate}
    \item \textbf{Numerical integration}
    \item \textbf{Monte Carlo simulation:}
    \begin{itemize}
        \item Sample $p_A^{(1)}, \ldots, p_A^{(S)}$ from $\text{Beta}(\alpha_A + x_A, \beta_A + n_A - x_A)$
        \item Sample $p_B^{(1)}, \ldots, p_B^{(S)}$ from $\text{Beta}(\alpha_B + x_B, \beta_B + n_B - x_B)$
        \item Estimate $P(p_B > p_A | D) \approx \frac{1}{S}\sum_{s=1}^S \mathbb{I}(p_B^{(s)} > p_A^{(s)})$
    \end{itemize}
\end{enumerate}

\begin{example}[Bayesian A/B Test]
\begin{align}
\text{Variant A: } & 80 \text{ conversions from } 1000 \text{ users} \\
\text{Variant B: } & 95 \text{ conversions from } 1000 \text{ users}
\end{align}

Using uniform priors $\text{Beta}(1, 1)$:
\begin{align}
p_A | D &\sim \text{Beta}(81, 921) \\
p_B | D &\sim \text{Beta}(96, 906)
\end{align}

Via Monte Carlo: $P(p_B > p_A | D) \approx 0.92$

Interpretation: There is a 92\% probability that variant B has higher conversion rate than A.
\end{example}

\lstset{language=R, caption=Bayesian A/B Test in R}
\begin{lstlisting}
# Bayesian A/B test comparison
bayesian_ab_test <- function(x_a, n_a, x_b, n_b,
                             prior_alpha = 1, prior_beta = 1,
                             n_simulations = 100000) {
  # Posterior parameters
  post_alpha_a <- prior_alpha + x_a
  post_beta_a <- prior_beta + n_a - x_a

  post_alpha_b <- prior_alpha + x_b
  post_beta_b <- prior_beta + n_b - x_b

  # Monte Carlo simulation
  p_a_samples <- rbeta(n_simulations, post_alpha_a, post_beta_a)
  p_b_samples <- rbeta(n_simulations, post_alpha_b, post_beta_b)

  # Probability B > A
  prob_b_better <- mean(p_b_samples > p_a_samples)

  # Expected lift
  lift <- p_b_samples / p_a_samples - 1
  expected_lift <- mean(lift)
  lift_ci <- quantile(lift, c(0.025, 0.975))

  # Difference distribution
  diff <- p_b_samples - p_a_samples
  diff_mean <- mean(diff)
  diff_ci <- quantile(diff, c(0.025, 0.975))

  list(
    prob_b_better = prob_b_better,
    expected_lift = expected_lift,
    lift_ci = lift_ci,
    difference_mean = diff_mean,
    difference_ci = diff_ci,
    posterior_a = list(alpha = post_alpha_a, beta = post_beta_a),
    posterior_b = list(alpha = post_alpha_b, beta = post_beta_b)
  )
}

# Example
result <- bayesian_ab_test(x_a = 80, n_a = 1000, x_b = 95, n_b = 1000)

cat("Probability B > A:", result$prob_b_better, "\n")
cat("Expected relative lift:", result$expected_lift, "\n")
cat("95% CI for lift:", result$lift_ci, "\n")
cat("Mean difference:", result$difference_mean, "\n")
cat("95% CI for difference:", result$difference_ci, "\n")
\end{lstlisting}

\subsection{Expected Loss}

Instead of binary decisions, we can quantify the expected loss from choosing the wrong variant.

\begin{definition}[Expected Loss]
The expected loss from choosing variant A is:
\begin{equation}
\mathcal{L}(A) = \mathbb{E}[\max(0, p_B - p_A) | D]
\end{equation}

Similarly for choosing B:
\begin{equation}
\mathcal{L}(B) = \mathbb{E}[\max(0, p_A - p_B) | D]
\end{equation}
\end{definition}

\textbf{Decision rule:} Choose the variant with smaller expected loss.

\textbf{Interpretation:} $\mathcal{L}(A)$ is the expected opportunity cost (in terms of conversion rate) if we choose A but B is actually better.

\lstset{language=Python, caption=Expected Loss Calculation}
\begin{lstlisting}
def expected_loss(x_a, n_a, x_b, n_b, prior_alpha=1, prior_beta=1,
                 n_simulations=100000):
    """
    Calculate expected loss for both variants.

    Returns:
    --------
    dict : Expected losses and recommendation
    """
    # Posterior parameters
    post_alpha_a = prior_alpha + x_a
    post_beta_a = prior_beta + n_a - x_a
    post_alpha_b = prior_alpha + x_b
    post_beta_b = prior_beta + n_b - x_b

    # Simulate from posteriors
    p_a = np.random.beta(post_alpha_a, post_beta_a, n_simulations)
    p_b = np.random.beta(post_alpha_b, post_beta_b, n_simulations)

    # Expected loss calculations
    loss_a = np.mean(np.maximum(0, p_b - p_a))
    loss_b = np.mean(np.maximum(0, p_a - p_b))

    # Relative expected loss
    rel_loss_a = loss_a / np.mean(p_a)
    rel_loss_b = loss_b / np.mean(p_b)

    recommendation = 'A' if loss_a < loss_b else 'B'

    return {
        'loss_a': loss_a,
        'loss_b': loss_b,
        'relative_loss_a': rel_loss_a,
        'relative_loss_b': rel_loss_b,
        'recommendation': recommendation,
        'loss_ratio': min(loss_a, loss_b) / max(loss_a, loss_b)
    }

# Example
result = expected_loss(x_a=80, n_a=1000, x_b=95, n_b=1000)

print(f"Expected loss choosing A: {result['loss_a']:.6f}")
print(f"Expected loss choosing B: {result['loss_b']:.6f}")
print(f"Relative loss A: {result['relative_loss_a']:.2%}")
print(f"Relative loss B: {result['relative_loss_b']:.2%}")
print(f"Recommendation: Variant {result['recommendation']}")
\end{lstlisting}

\section{Prior Selection}

Prior choice affects posterior, especially with limited data.

\subsection{Types of Priors}

\begin{enumerate}
    \item \textbf{Uninformative (Flat) Priors}

    $\text{Beta}(1, 1)$ (uniform) or Jeffrey's prior $\text{Beta}(0.5, 0.5)$

    Lets data dominate; minimal assumptions

    \item \textbf{Weakly Informative Priors}

    $\text{Beta}(2, 20)$ for expected conversion around 10\%

    Provides gentle regularization without strong assumptions

    \item \textbf{Informative Priors}

    Based on historical data or expert knowledge

    Example: $\text{Beta}(50, 450)$ from previous experiments showing 10\% conversion

    \item \textbf{Empirical Bayes}

    Use data itself to estimate prior (controversial but practical)
\end{enumerate}

\subsection{Prior Sensitivity Analysis}

Always check how conclusions depend on prior choice.

\begin{example}[Prior Sensitivity]
Test with 10 conversions from 100 users:

\begin{center}
\begin{tabular}{lccc}
\toprule
\textbf{Prior} & \textbf{Post. Mean} & \textbf{Post. Std} & \textbf{95\% CI} \\
\midrule
$\text{Beta}(1, 1)$ & 0.108 & 0.031 & [0.058, 0.175] \\
$\text{Beta}(5, 45)$ & 0.106 & 0.026 & [0.063, 0.162] \\
$\text{Beta}(10, 90)$ & 0.105 & 0.023 & [0.066, 0.154] \\
$\text{Beta}(20, 180)$ & 0.103 & 0.020 & [0.069, 0.145] \\
\bottomrule
\end{tabular}
\end{center}

With small data, stronger priors pull estimates toward prior mean.
\end{example}

\section{Advantages of Bayesian A/B Testing}

\begin{enumerate}
    \item \textbf{Direct Probability Statements}

    "95\% probability that B is better" vs. "reject null with p = 0.03"

    More intuitive for stakeholders

    \item \textbf{Natural Sequential Testing}

    Can stop anytime without inflating error rates

    No need for $\alpha$-spending functions

    \item \textbf{Incorporation of Prior Information}

    Leverage historical data and expert knowledge

    Particularly valuable for rare events

    \item \textbf{Coherent Decision Framework}

    Expected loss provides principled stopping rules

    Balances statistical and business considerations

    \item \textbf{Full Uncertainty Quantification}

    Entire posterior distribution available

    Can compute any quantity of interest
\end{enumerate}

\section{Limitations and Challenges}

\begin{enumerate}
    \item \textbf{Prior Dependence}

    Results depend on prior, especially with little data

    Requires prior specification and sensitivity analysis

    \item \textbf{Computational Complexity}

    Complex models may require MCMC or variational inference

    Not always analytical like conjugate examples

    \item \textbf{Interpretation Challenges}

    "Probability B > A" requires careful framing

    Can be misinterpreted as certainty

    \item \textbf{Less Familiar}

    Stakeholders more familiar with p-values and confidence intervals

    Requires education and buy-in

    \item \textbf{No Error Rate Guarantees}

    Unlike frequentist methods, no Type I/II error control

    Different philosophical framework for errors
\end{enumerate}

\section{Bayesian vs. Frequentist: When to Use Which}

\begin{center}
\begin{tabular}{lcc}
\toprule
\textbf{Scenario} & \textbf{Frequentist} & \textbf{Bayesian} \\
\midrule
Fixed sample size & \checkmark & \checkmark \\
Sequential/adaptive & & \checkmark \\
Prior information available & & \checkmark \\
Regulatory requirement & \checkmark & \\
Simple communication & \checkmark & \\
Probability statements needed & & \checkmark \\
Multiple metrics/variants & \checkmark & \checkmark \\
\bottomrule
\end{tabular}
\end{center}

\textbf{Recommendation:} Both frameworks are valuable. Choose based on:
\begin{itemize}
    \item Organizational culture and stakeholder preferences
    \item Regulatory requirements
    \item Prior information availability
    \item Sequential testing needs
    \item Computational resources
\end{itemize}

Many organizations use hybrid approaches: frequentist for primary analysis, Bayesian for exploration.

\section{Summary}

Bayesian A/B testing offers a powerful alternative to frequentist methods:

\begin{itemize}
    \item Treats parameters as random variables with probability distributions
    \item Updates beliefs via Bayes' theorem
    \item Conjugate priors enable analytical solutions
    \item Provides direct probability statements
    \item Handles sequential testing naturally
    \item Requires prior specification and sensitivity analysis
\end{itemize}

The choice between Bayesian and frequentist approaches depends on context, but understanding both enriches the practitioner's toolkit.

\section{Further Reading}

\begin{itemize}
    \item \cite{gelman2013bayesian}: Comprehensive Bayesian data analysis
    \item \cite{kruschke2014doing}: Bayesian analysis for psychology (accessible)
    \item \cite{stucchio2015bayesian}: Bayesian A/B testing for practitioners
\end{itemize}

\section{Exercises}

\begin{enumerate}
    \item \textbf{Beta-Binomial Calculations}
    \begin{enumerate}
        \item Derive the posterior for Beta-Binomial model from first principles.
        \item Show that $\text{Beta}(1,1)$ is the uniform distribution on $[0,1]$.
    \end{enumerate}

    \item \textbf{Bayesian A/B Test}
    \begin{enumerate}
        \item Implement complete Bayesian A/B test including posterior visualization, probability calculations, and expected loss.
        \item Compare Bayesian and frequentist conclusions for various sample sizes.
    \end{enumerate}

    \item \textbf{Prior Sensitivity}
    \begin{enumerate}
        \item Conduct sensitivity analysis varying prior from uninformative to strongly informative.
        \item At what sample size do priors matter little?
    \end{enumerate}

    \item \textbf{Sequential Testing}
    \begin{enumerate}
        \item Simulate sequential Bayesian testing with stopping rule based on probability threshold.
        \item Compare to frequentist sequential testing in terms of sample size and error rates.
    \end{enumerate}
\end{enumerate}
